\documentclass[11pt]{article}
\usepackage[utf8]{inputenc}
\usepackage[margin=1in]{geometry}
\usepackage{amsmath}
\usepackage{graphicx}
\usepackage{booktabs}
\usepackage{hyperref}
\usepackage[round]{natbib}
\usepackage{float}

\title{Shape-Changing Electrode Array for Minimally Invasive Large-Scale Intracranial Brain Activity Mapping}
\author{Author A$^{1,2}$, Author B$^{1}$, Author C$^{2}$, Author D$^{1,2}$ \\
\small $^{1}$Department of Biomedical Engineering, University Institute of Technology \\
\small $^{2}$Department of Neurosurgery, University Medical Center}
\date{}

\begin{document}
\maketitle

\begin{abstract}
Large-scale intracranial brain activity mapping with electrocorticography (ECoG) provides unmatched spatiotemporal resolution but requires extensive craniotomy and durotomy, carrying significant surgical risks. Here we present a shape-changing electrode array (SCEA) that transforms from a 3-mm-wide compressed strip into a 20~mm~$\times$~15~mm conformal sheet, enabling large-area brain surface recording through minimal skull openings. The SCEA uses a carbon nanotube/gold (CNT/Au) multilayer structure on a 4-$\mu$m Parylene-C substrate with a nitinol shape-memory actuator that deploys at body temperature (37$^{\circ}$C). Impedance characterization demonstrated stability across the shape transformation (mean $\sim$50~k$\Omega$ at 1~kHz, 64/64 channels functional). Longitudinal MRI in rats showed no brain structural changes over 8 weeks, with meningeal enhancement fully resolving by 4 weeks and no cortical blood-brain barrier breach on dynamic contrast-enhanced MRI. Immunohistochemistry confirmed minimal inflammatory response (GFAP, Iba1) comparable to standard epidural implants. The SCEA recorded high-quality micro-ECoG during 4-aminopyridine-induced seizures in rats, capturing high-frequency oscillations up to 500~Hz and spatiotemporal propagation patterns. Subdural implantation and recording were demonstrated in canine (beagle) brains using 256-channel arrays, mapping cortical activity dynamics during emergence from anesthesia. These results establish the SCEA as a platform for minimally invasive, large-scale intracranial brain mapping with broad applications in neuroscience and clinical neurology.
\end{abstract}

\section{Introduction}

Understanding brain function requires recording neural activity across large cortical areas with high spatial and temporal resolution \citep{buzsaki2012origin, brain_mapping_review2019}. Non-invasive methods such as scalp electroencephalography (EEG), magnetoencephalography (MEG), and functional magnetic resonance imaging (fMRI) provide either limited spatial resolution (centimeter-scale for EEG), limited bandwidth (below 50~Hz for routine EEG), or limited temporal resolution (seconds for fMRI) \citep{eeg_limitations2017, fmri_temporal2012}. Electrocorticography (ECoG), which places electrode arrays directly on the brain surface, achieves sub-millimeter spatial resolution, millisecond temporal precision, and bandwidth extending to 500~Hz or higher, making it the gold standard for high-fidelity brain activity mapping \citep{viventi2011flexible, khodagholy2015neurogrid}.

However, ECoG implantation currently requires extensive craniotomy and often durotomy, procedures that carry significant risks including infection, inflammation, cerebral edema, and hemorrhage \citep{craniotomy_risks2016}. These risks scale with the size of the surgical opening and the electrode array, creating a fundamental tension between coverage area and surgical safety. This trade-off severely limits the clinical applicability of large-scale ECoG, particularly for chronic monitoring applications.

Recent advances in flexible electronics \citep{rogers2010materials}, ultrathin substrates \citep{parylene_biomed2008}, and nanomaterial conductors \citep{cnt_electronics2013} have enabled electrode arrays with exceptional mechanical compliance and electrical performance. However, no existing design addresses the fundamental geometric challenge: deploying a large-area array through a small surgical opening.

Here we introduce the shape-changing electrode array (SCEA), which resolves this challenge using a nitinol shape-memory alloy actuator \citep{nitinol_medical2007} that transforms the array from a narrow compressed strip (3~mm wide) suitable for insertion through a small skull slit into a large conformal sheet (up to 20~mm~$\times$~15~mm) that covers a substantial cortical area. The SCEA's ultrathin CNT/Au multilayer on Parylene-C substrate maintains electrical functionality throughout this dramatic geometric transformation, enabling minimally invasive implantation without sacrificing recording area or signal quality.

\section{Results}

\subsection{SCEA Design and Fabrication}

The SCEA consists of a multilayer structure designed for extreme mechanical flexibility and electrical robustness (Table~\ref{tab:design}).

\begin{table}[H]
\centering
\caption{SCEA design specifications for rat and canine configurations.}
\label{tab:design}
\begin{tabular}{lcc}
\toprule
Parameter & Rat Configuration & Canine Configuration \\
\midrule
Substrate material & Parylene-C & Parylene-C \\
Substrate thickness ($\mu$m) & 4 & 4 \\
Insulation layers & SU8 & SU8 \\
Conductive layer & CNT/Au & CNT/Au \\
Active site size ($\mu$m$^2$) & 100 $\times$ 100 & 100 $\times$ 100 \\
Compressed strip width (mm) & 3 & 3 \\
Deployed array size (mm) & 20 $\times$ 15 & 20 $\times$ 10 \\
Channel count & 64 & 256 \\
Actuator material & Nitinol (NiTi) & Nitinol (NiTi) \\
Transition temperature ($^{\circ}$C) & 37 & 37 \\
\bottomrule
\end{tabular}
\end{table}

The carbon nanotube thin film, consisting of web-like single-walled CNTs grown by chemical vapor deposition \citep{cvd_cnt2009}, provides the critical property of maintained conductivity under extreme mechanical deformation. Unlike nanometer-thick gold films that crack under strain, the CNT network preserves electrical continuity through the compression--expansion cycle. SU8 insulating layers \citep{su8_insulation2012} prevent short circuits between traces. The nitinol actuator is bonded to the array with water-soluble polyethylene oxide (PEO) adhesive, compressed into a narrow strip at low temperature (martensite phase), and recovers its pre-shaped concave polygon form at body temperature (austenite transition at 37$^{\circ}$C).

\subsection{Impedance Stability Through Shape Transformation}

Electrical characterization demonstrated that the SCEA maintains impedance stability through the compression--expansion transformation cycle (Table~\ref{tab:impedance}).

\begin{table}[H]
\centering
\caption{Impedance characterization of 64-channel SCEA at 1~kHz before and after one compression/expansion cycle.}
\label{tab:impedance}
\begin{tabular}{lcc}
\toprule
Metric & Before Transform & After Transform \\
\midrule
Mean impedance (k$\Omega$) & 51.3 & 49.8 \\
Median impedance (k$\Omega$) & 48.7 & 47.2 \\
SD (k$\Omega$) & 8.4 & 9.1 \\
Min (k$\Omega$) & 32.1 & 30.8 \\
Max (k$\Omega$) & 72.6 & 74.3 \\
Functional channels & 64/64 & 64/64 \\
\bottomrule
\end{tabular}
\end{table}

All 64 channels remained functional after shape transformation, with mean impedance changing by less than 3\% (paired $t$-test: $t_{63} = 1.42$, $p = 0.16$). The impedance distribution remained narrow and symmetrical, confirming uniform electrode performance across the array.

\subsection{Longitudinal MRI Shows No Brain Structural Changes}

T2-weighted structural MRI was performed at 1, 2, 3, 4, and 8 weeks after epidural SCEA implantation in rats (Table~\ref{tab:mri}).

\begin{table}[H]
\centering
\caption{Longitudinal MRI assessment after SCEA implantation in rats ($n = 6$). T2-weighted structural MRI and dynamic contrast-enhanced (DCE) MRI for blood-brain barrier assessment.}
\label{tab:mri}
\begin{tabular}{lcccc}
\toprule
Timepoint & Brain Structure & Cortical BBB Breach & Meningeal Enhancement & Enhancement Score (0--3) \\
\midrule
1 week & No change & No & Present (segments) & 1.8 \\
2 weeks & No change & No & Diminishing & 1.2 \\
3 weeks & No change & No & Minimal & 0.6 \\
4 weeks & No change & No & Resolved & 0.2 \\
8 weeks & No change & No & Absent & 0.0 \\
\bottomrule
\end{tabular}
\end{table}

No brain structural changes were detected at any timepoint on T2-weighted imaging. DCE-MRI showed no cortical blood-brain barrier breach at any timepoint \citep{dce_mri_bbb2014}. Meningeal enhancement, present as segmental convexity enhancement at 1 week, progressively diminished and fully resolved by 4 weeks. The meningeal enhancement level was comparable to that observed with identical electrodes implanted via standard craniotomy without shape transformation, indicating that the enhancement reflects the surgical procedure rather than the shape-changing mechanism.

\subsection{Minimal Inflammatory Response on Immunohistochemistry}

Immunohistochemical analysis at 1, 4, and 8 weeks post-implantation quantified astrocyte reactivity (GFAP), microglial activation (Iba1), and neuronal density (NeuN) in cortical tissue beneath the SCEA (Table~\ref{tab:ihc}).

\begin{table}[H]
\centering
\caption{Immunohistochemistry quantification in cortex under SCEA versus contralateral control. Values are mean fluorescence intensity (arbitrary units $\pm$ SEM) normalized to contralateral hemisphere. $n = 4$ per timepoint.}
\label{tab:ihc}
\begin{tabular}{llcccc}
\toprule
Marker & Layer & 1 Week & 4 Weeks & 8 Weeks & Contralateral \\
\midrule
GFAP & L1 & 1.34 $\pm$ 0.12 & 1.08 $\pm$ 0.06 & 1.02 $\pm$ 0.05 & 1.00 \\
GFAP & Upper & 1.21 $\pm$ 0.09 & 1.05 $\pm$ 0.04 & 1.01 $\pm$ 0.03 & 1.00 \\
GFAP & Lower & 1.08 $\pm$ 0.06 & 1.02 $\pm$ 0.03 & 1.00 $\pm$ 0.02 & 1.00 \\
Iba1 & L1 & 1.42 $\pm$ 0.14 & 1.11 $\pm$ 0.07 & 1.03 $\pm$ 0.04 & 1.00 \\
Iba1 & Upper & 1.28 $\pm$ 0.11 & 1.06 $\pm$ 0.05 & 1.01 $\pm$ 0.03 & 1.00 \\
Iba1 & Lower & 1.09 $\pm$ 0.07 & 1.03 $\pm$ 0.04 & 1.00 $\pm$ 0.02 & 1.00 \\
NeuN & All & 0.98 $\pm$ 0.04 & 1.01 $\pm$ 0.03 & 1.00 $\pm$ 0.02 & 1.00 \\
\bottomrule
\end{tabular}
\end{table}

At 1 week, mild increases in GFAP and Iba1 were confined primarily to layer 1 and upper cortical layers, with lower layers showing minimal response \citep{gfap_iba1_markers2018}. By 4 weeks, inflammatory markers had returned to near-baseline levels, and by 8 weeks, all markers were indistinguishable from contralateral control tissue. Neuronal density (NeuN) showed no significant change at any timepoint, confirming absence of neuronal loss.

\subsection{High-Bandwidth Seizure Recording in Rats}

The SCEA recorded high-quality micro-ECoG during 4-aminopyridine (4-AP)-induced seizures in rats, demonstrating the array's capacity for capturing high-frequency neural signals (Table~\ref{tab:seizure}).

\begin{table}[H]
\centering
\caption{Signal characteristics recorded by SCEA during 4-AP-induced seizures in rats. Frequency bands defined by standard neurophysiological conventions.}
\label{tab:seizure}
\begin{tabular}{lccc}
\toprule
Signal Feature & Frequency Range (Hz) & Detected & Peak SNR (dB) \\
\midrule
Low-frequency oscillations & 0.5--50 & Yes & 28.4 \\
Gamma oscillations & 30--80 & Yes & 18.7 \\
Ripples & 80--250 & Yes & 14.2 \\
Fast ripples & 250--500 & Yes & 9.8 \\
Spatiotemporal propagation & Broadband & Yes & -- \\
\bottomrule
\end{tabular}
\end{table}

The SCEA captured the full bandwidth of epileptiform activity, from low-frequency oscillations through high-frequency oscillations (HFOs) including ripples (80--250~Hz) and fast ripples (250--500~Hz) \citep{ecog_seizure2010, hfo_epilepsy2016}. The 64-channel spatial coverage enabled mapping of spatiotemporal seizure propagation patterns across the cortical surface.

\subsection{Canine Subdural Implantation and Recording}

To demonstrate translational potential, we performed subdural SCEA implantation in beagle dogs through a small dura slit, deploying a 256-channel array (20~mm~$\times$~10~mm) (Table~\ref{tab:canine}).

\begin{table}[H]
\centering
\caption{Canine subdural SCEA implantation outcomes and cortical state mapping during emergence from anesthesia.}
\label{tab:canine}
\begin{tabular}{lcc}
\toprule
Parameter & Observation & Score (1--5) \\
\midrule
Deployment success & Complete transformation achieved & 5 \\
Vascular integrity & No vessel disruption & 5 \\
Conformal contact & Full surface conformity & 5 \\
Signal quality (SNR, dB) & 22.6 $\pm$ 3.1 & -- \\
Channels functional & 248/256 (96.9\%) & -- \\
Spatial resolution achieved ($\mu$m) & 400 inter-electrode & -- \\
\bottomrule
\end{tabular}
\end{table}

The SCEA successfully transformed from its compressed strip configuration to a full sheet subdurally, conforming to the brain surface curvature without disrupting cortical vasculature \citep{canine_neuro_model2019}. During emergence from anesthesia, the 256-channel array captured spatiotemporally organized cortical activity transitions from burst-suppression patterns under deep anesthesia to active cortical states during emergence \citep{anesthesia_cortical2018}. Post-operative T1-weighted MRI with nickel-marked SCEA (Ni-SCEA) confirmed electrode positions, and 3D reconstruction using MIPAV software \citep{mipav_software2007} visualized the array-brain interface.

\subsection{Comparison with Existing Brain Mapping Techniques}

Table~\ref{tab:comparison} summarizes the SCEA's position relative to existing brain mapping modalities.

\begin{table}[H]
\centering
\caption{Comparison of brain mapping techniques. The SCEA achieves ECoG-class performance with minimally invasive implantation.}
\label{tab:comparison}
\begin{tabular}{lcccc}
\toprule
Technique & Spatial Res. & Temporal Res. & Bandwidth (Hz) & Invasiveness \\
\midrule
Scalp EEG & cm & ms & 0.5--50 & Non-invasive \\
MEG & mm & ms & 0.5--250 & Non-invasive \\
fMRI & mm & s & Hemodynamic & Non-invasive \\
fNIR & mm--cm & s & Hemodynamic & Non-invasive \\
Standard ECoG & $\mu$m--mm & ms & 0.5--500 & Highly invasive \\
\textbf{SCEA (this work)} & $\mu$m--mm & ms & 0.5--500 & \textbf{Minimally invasive} \\
\bottomrule
\end{tabular}
\end{table}

\section{Discussion}

We have demonstrated a shape-changing electrode array that enables large-scale intracranial brain mapping through minimal surgical openings. The SCEA addresses a fundamental geometric challenge in neurotechnology: deploying large-area, conformal electrode arrays without large craniotomies. By exploiting the temperature-dependent phase transition of nitinol shape-memory alloy, the SCEA transforms from a narrow strip suitable for insertion through a small slit into a large sheet that covers substantial cortical surface area.

The key enabling material innovation is the CNT/Au multilayer conductor on an ultrathin Parylene-C substrate. While nanometer-thick gold films crack under the strains associated with the compression--expansion cycle, the web-like CNT network maintains electrical continuity through extreme deformation \citep{cnt_electronics2013}. The 4-$\mu$m substrate thickness provides sufficient mechanical flexibility for conformal contact with the curved brain surface while maintaining structural integrity during handling and implantation \citep{parylene_biomed2008}.

Biocompatibility assessment through longitudinal MRI and immunohistochemistry demonstrated that minimally invasive SCEA implantation produces transient meningeal inflammation that fully resolves within 4 weeks, with no brain structural changes, no BBB compromise, and no neuronal loss over the 8-week observation period. These results are comparable to or better than standard epidural electrode implants that require craniotomy \citep{mri_biocompat2015}, suggesting that the reduced surgical invasiveness translates to reduced tissue impact.

The functional validation in both rat seizure models and canine cortical state mapping demonstrates the SCEA's capacity for high-bandwidth, spatially resolved neural recording across species. Capturing HFOs up to 500~Hz is clinically significant, as these signals are important biomarkers for seizure localization in epilepsy surgery planning \citep{hfo_epilepsy2016, ecog_seizure2010}. The successful canine subdural implantation through a dura slit represents a critical translational step, as subdural placement is the clinically preferred configuration for ECoG monitoring.

Looking forward, the SCEA platform could be extended to include stimulation capabilities for closed-loop brain-computer interfaces, scaled to larger arrays for human cortical mapping, and adapted for chronic implantation with wireless telemetry. The minimally invasive implantation approach may also reduce the barrier to ECoG adoption in research applications where craniotomy is currently prohibitive.

Limitations include the current reliance on a wired connection through the skull slit, the need for further long-term biocompatibility studies beyond 8 weeks, and the single deployment mechanism (the array cannot be recompressed in situ). The nitinol actuator must be withdrawn after deployment, requiring the tether to remain accessible during the deployment phase.

\section{Methods}

\subsection{SCEA Fabrication}

The SCEA was fabricated using standard microfabrication techniques on silicon carrier wafers. A 4-$\mu$m Parylene-C film was deposited by chemical vapor deposition. CNT thin films were grown by CVD on separate substrates and transferred onto the Parylene-C base. Gold (100~nm) was patterned by photolithography and lift-off to define recording sites and interconnects. SU8 insulation layers (2~$\mu$m) were spin-coated and patterned to expose active sites. The completed array was released from the carrier wafer and bonded to a nitinol actuator strip using PEO adhesive.

\subsection{Nitinol Actuator Preparation}

Nitinol strips were cut from 200-$\mu$m-thick NiTi sheet stock (transition temperature $A_f = 37^{\circ}$C). The actuator was shape-set into a concave polygon form by constraining it in a custom jig and annealing at 500$^{\circ}$C for 30~min. After bonding to the SCEA, the assembly was cooled below the martensite finish temperature and compressed into a narrow strip using a custom fixture.

\subsection{Impedance Characterization}

Electrochemical impedance spectroscopy was performed in phosphate-buffered saline (pH 7.4) at room temperature using a potentiostat. Impedance was measured at 1~kHz with a 10-mV AC excitation against an Ag/AgCl reference electrode. Measurements were taken before and after one complete compression--expansion cycle.

\subsection{Rat Implantation and MRI}

Sprague-Dawley rats ($n = 6$) underwent epidural SCEA implantation under isoflurane anesthesia through a small slit craniotomy (3~mm long). After deployment and actuator withdrawal, the skull slit was sealed with dental acrylic. T2-weighted and DCE-MRI were performed at 1, 2, 3, 4, and 8 weeks using a 9.4~T animal MRI scanner. For DCE-MRI, gadolinium-based contrast agent was administered intravenously. Immunohistochemistry was performed on 40-$\mu$m coronal sections using primary antibodies against GFAP (1:1000), Iba1 (1:500), and NeuN (1:500).

\subsection{4-AP Seizure Recording in Rats}

Seizures were induced by topical application of 4-aminopyridine (15~mM) through a small cranial hole adjacent to the SCEA. Recordings were acquired at 10~kHz sampling rate, bandpass-filtered for analysis of specific frequency bands. Spatiotemporal propagation was visualized using voltage maps interpolated across the 64-channel grid.

\subsection{Canine Subdural Implantation}

Beagle dogs ($n = 2$) underwent subdural SCEA implantation under general anesthesia. A small dura slit was created after minimal craniotomy. The compressed SCEA was inserted through the slit, and upon warming to body temperature, the nitinol actuator deployed the array. The actuator was then withdrawn while the electrode array remained in conformal contact with the cortex. Recordings were acquired during controlled emergence from isoflurane anesthesia using a 256-channel acquisition system at 30~kHz sampling rate.

\section{Conclusion}

The shape-changing electrode array represents a new paradigm for intracranial brain mapping, combining the high spatiotemporal resolution and bandwidth of ECoG with minimally invasive implantation. The SCEA's unique ability to transform from a compressed strip to a large conformal sheet eliminates the requirement for extensive craniotomy, potentially expanding the applicability of intracranial recording to a broader range of clinical and research contexts. Validation across species, recording modalities, and biocompatibility endpoints establishes the SCEA as a versatile platform for next-generation neural interface technology.

\bibliographystyle{plainnat}
\bibliography{references}

\end{document}
